\documentclass[11pt]{article}
\usepackage[T1]{fontenc}
\usepackage[utf8]{inputenc}
\usepackage[a4paper,margin=1in]{geometry}
\usepackage{amsmath,amssymb,amsthm,mathtools}
% microtype not used in this environment
\usepackage{xcolor}
\usepackage{hyperref}
\usepackage{enumitem}
\usepackage{setspace}
% lmodern not available in this environment
\usepackage{booktabs}
\usepackage{array}
\hypersetup{colorlinks=true,linkcolor=blue!50!black,
urlcolor=blue!50!black,citecolor=blue!50!black}
\setlength{\parskip}{0.5em}
\setlength{\parindent}{0pt}
\onehalfspacing

\title{\vspace{-1.5em}\bfseries
From Local Curvature to the Weil Functional:\\[0.4em]
\large An Explicit Construction}
\author{Ulrich Tehrani}
% FIX v5 — version update
\date{Research Note \textperiodcentered\ April 2026 \textperiodcentered\ v5}

% v5 corrections (April 2026):
% (1) FIX v5: Introduction framing softened — no longer implies
%     that the paper "packages" a global Weil positivity statement;
%     only the prime-side construction is claimed.
% (2) FIX v5: Abstract clarified — spectral inequality marked as
%     open and of RH strength.
% (3) FIX v5: Remark[Status of the spectral inequality] added
%     before eq:openineq with explicit open/RH-strength statement.
%
% v4 corrections (March 2026):
% (1) V_p(1/2) table values corrected (3.844, 3.621, 3.238, 2.945, 2.530)
% (2) c_p^ren simplified to f_p^{1/2} (removed p^{1/4}/(log p)^{1/2} factor)
%     so that [c_p^ren]^2 = f_p exactly, consistent with D = sum f_p
% (3) H_ren asymptotic corrected: H_ren has sawtooth structure,
%     not simple log kappa growth; Mertens connection clarified

\newtheorem*{lemma*}{Lemma}
\newtheorem*{proposition*}{Proposition}
\newtheorem*{corollary*}{Corollary}
\newtheorem*{remark*}{Remark}

\begin{document}
\maketitle

\begin{abstract}
Building on the curvature decomposition established in~\cite{paper1},
we construct an explicit family of admissible test functions
$g^*_{\sigma,\varepsilon} \in S_{ad}$ for the Weil explicit formula.
For this family we derive a prime-side identity of the form
\[
W(g^*_{\sigma,\varepsilon} * \check g^*_{\sigma,\varepsilon})
=
Z(g^*_{\sigma,\varepsilon}) - H_{\mathrm{local}}(\sigma,\kappa) + O(\varepsilon),
\]
where $Z(g^*_{\sigma,\varepsilon})$ denotes the corresponding zero sum.
We further introduce renormalized prime weights
$c_p^{\mathrm{ren}} > 0$, prove convergence of the associated series
\[
D = \sum_p \bigl(c_p^{\mathrm{ren}}\bigr)^2 \approx 9.471,
\]
and show that the renormalized prime-side scale is logarithmic in $\kappa$.
% FIX v5 — clarify RH-strength of open inequality
This isolates the remaining open problem as a single spectral inequality
relating $Z(g^*_{1/2,\varepsilon})$ to $H_{\mathrm{local}}(1/2,\kappa)$;
establishing this inequality unconditionally would be of RH strength.
We also record the appearance of the Mertens constant as the natural scale
of the renormalized curvature remainder.
\end{abstract}

\section{Introduction and relation to Paper~1}

In~\cite{paper1} we studied the curvature field
\[
H_{\xi}(\sigma,t)
:= \frac{d^2}{d\sigma^2}\log|\xi(\sigma+it)|^2
= H_{\mathrm{local}}(\sigma,\kappa) + H_{\mathrm{dual}}(\sigma,t,\kappa),
\]
where $H_{\mathrm{dual}} := H_{\xi} - H_{\mathrm{local}}$ is the spectral
remainder encoding the zero contributions (sign convention as in~\cite{paper1},
equation~(5)).
where the truncated local curvature
\[
H_{\mathrm{local}}(\sigma,\kappa) = \psi_1(\sigma)
+ \sum_{p\le\kappa} V_p(\sigma)
\]
contains the archimedean term and the finite-place contributions. The prime-side pieces satisfy
$V_p(\sigma) \ge 0$, and on the critical line one has
$H_{\mathrm{local}}(1/2,\kappa) \asymp (\log\kappa)^2$.

% FIX v5 — soften framing
The central open question of~\cite[\S7]{paper1} asks whether this local
positivity and critical divergence can be organized into an explicit admissible
test-function construction whose prime-side term matches the local curvature;
the spectral (zero-side) inequality remains open. The present note addresses a
concrete part of that question by constructing an explicit family
$g^*_{\sigma,\varepsilon} \in S_{ad}$ and identifying the corresponding
prime-side term in the Weil functional with the local curvature studied in
Paper~1.

The remaining open step is then isolated as a single spectral inequality;
see Section~\ref{sec:open}.

\section{Construction of the admissible test function}

\subsection{The raw family}

Recall from~\cite[Eq.~(3)]{paper1} that
\[
V_p(\sigma)
= 4(\log p)^2\,\frac{p^{-2\sigma}}{(1-p^{-2\sigma})^2}.
\]
Define
\[
c_{p,\sigma}
:= V_p(\sigma)^{1/2}
\cdot \frac{p^{1/4}}{(\log p)^{1/2}},
\]
and let
\[
\varphi_\varepsilon(x)
= \frac{1}{\varepsilon\sqrt{2\pi}}e^{-x^2/(2\varepsilon^2)}
\]
be the standard Gaussian bump. For fixed $\kappa$ set
\[
\widehat f_{\sigma,\varepsilon}(x)
= \sum_{p\le\kappa} c_{p,\sigma}\,\varphi_\varepsilon(x-\log p).
\]
The admissible test function is obtained by antisymmetrization:
\begin{equation}\label{eq:gstar}
\widehat g^*_{\sigma,\varepsilon}(x)
:= \widehat f_{\sigma,\varepsilon}(x)
- \widehat f_{\sigma,\varepsilon}(-x).
\end{equation}

\begin{lemma*}[$g^*_{\sigma,\varepsilon} \in S_{ad}$]
For every $\sigma>0$ and $\varepsilon>0$, the function
$g^*_{\sigma,\varepsilon}$ belongs to the admissible class $S_{ad}$ used in
Lagarias' formulation of Weil positivity.
\end{lemma*}

\textit{Proof.}
The Gaussian bumps are Schwartz, and antisymmetrization preserves the
Schwartz class. Hence $g^*_{\sigma,\varepsilon}$ and
$\widehat g^*_{\sigma,\varepsilon}$ are rapidly decreasing. Moreover,
$\widehat g^*_{\sigma,\varepsilon}(0)=0$ by antisymmetry, and likewise
$g^*_{\sigma,\varepsilon}(0)=0$. These are the required vanishing conditions
at the origin. Finally, the Gaussian decay ensures absolute convergence of the
prime-side and zero-side expressions appearing in the Weil functional for this
family. \hfill$\square$

\section{Prime-side identity for the Weil functional}

\begin{proposition*}[Prime-side identity]
For the family $g^*_{\sigma,\varepsilon}$ constructed above, one has
\begin{equation}\label{eq:weil-decomp}
W(g^*_{\sigma,\varepsilon} * \check g^*_{\sigma,\varepsilon})
=
Z(g^*_{\sigma,\varepsilon}) - H_{\mathrm{local}}(\sigma,\kappa) + O(\varepsilon),
\end{equation}
where
$Z(g^*_{\sigma,\varepsilon}) = \sum_{\rho}
|\widehat g^*_{\sigma,\varepsilon}(\rho)|^2$
runs over the nontrivial zeros of $\zeta(s)$.
\end{proposition*}

\textit{Proof sketch.}
Apply the Weil explicit formula in the Lagarias framework. The constant
term vanishes because $\widehat g^*_{\sigma,\varepsilon}(0)=0$.
By construction of the coefficients $c_{p,\sigma}$, the prime-side quadratic
form converges to $H_{\mathrm{local}}(\sigma,\kappa)$ as $\varepsilon\to 0$.
The error term $O(\varepsilon)$ arises from the Gaussian smoothing of the
prime-power contributions and is not made uniform in $\kappa$ in the present
note. \hfill$\square$

\begin{remark*}
Equation~\eqref{eq:weil-decomp} is the key bridge between the curvature field
of Paper~1 and the Weil positivity framework. The archimedean
term $\hat{g}^*(0)^2 c_0$ vanishes since $\hat{g}^*(0)=0$
by condition~(S2), so the archimedean contribution to~$W$
is zero in the present normalization.
\end{remark*}

\section{Renormalized weights and convergence}

\subsection{Positivity of the renormalized coefficients}

At $\sigma=1/2$ define
\[
f_p := V_p(1/2) - \frac{4(\log p)^2}{p}.
\]
A direct calculation gives
\begin{equation}\label{eq:fp}
f_p
= \frac{4(\log p)^2(2p-1)}{p(p-1)^2} > 0
\qquad (p\ge 2).
\end{equation}

\begin{lemma*}[Renormalized coefficients are positive]
For every prime $p\ge 2$, the quantity $f_p$ in~\eqref{eq:fp}
is strictly positive.
\end{lemma*}

\textit{Proof.}
The numerator $2p-1>0$ and the denominator $p(p-1)^2>0$
for every prime $p\ge 2$. \hfill$\square$

The renormalized coefficients are simply
\begin{equation}\label{eq:cren}
c_p^{\mathrm{ren}} := f_p^{1/2} > 0,
\end{equation}
so that $[c_p^{\mathrm{ren}}]^2 = f_p$.
Numerical values for small primes:

\begin{center}
\begin{tabular}{ccccc}
\toprule
$p$ & $V_p(1/2)$ & $4(\log p)^2/p$ & $f_p = [c_p^{\mathrm{ren}}]^2$ \\
\midrule
2  & 3.844 & 0.961 & 2.883 \\
3  & 3.621 & 1.609 & 2.012 \\
5  & 3.238 & 2.072 & 1.166 \\
7  & 2.945 & 2.164 & 0.781 \\
11 & 2.530 & 2.091 & 0.439 \\
\bottomrule
\end{tabular}
\end{center}

\subsection{Convergence of the renormalized energy}

\begin{lemma*}[Convergence]
The series
\begin{equation}\label{eq:D}
D := \sum_p \bigl[c_p^{\mathrm{ren}}\bigr]^2
   = \sum_p \frac{4(\log p)^2(2p-1)}{p(p-1)^2}
\end{equation}
converges absolutely. Numerically $D \approx 9.471$
(primes below $10^4$; tail $<0.009$).
\end{lemma*}

\textit{Proof.}
For large $p$, $[c_p^{\mathrm{ren}}]^2 \sim 8(\log p)^2/p^2$.
Since $\sum_n (\log n)^2/n^2 <\infty$, convergence follows
by comparison. \hfill$\square$

\section{The renormalized scale and the Mertens constant}

By the prime number theorem,
\[
H_{\mathrm{local}}(1/2,\kappa)
\sim 2(\log\kappa)^2
\qquad(\kappa\to\infty),
\]
so the renormalized remainder
\[
H_{\mathrm{local}}^{\mathrm{ren}}(1/2,\kappa)
:= H_{\mathrm{local}}(1/2,\kappa) - 2(\log\kappa)^2
\]
is of lower order than $(\log\kappa)^2$.
Numerically, $H_{\mathrm{local}}^{\mathrm{ren}}$ exhibits
a sawtooth pattern: it jumps upward at each new prime $p\le\kappa$
by $f_p > 0$, and decreases continuously between primes as
$2(\log\kappa)^2$ grows. This is the direct analogue of
the prime-counting function $\pi(x)$, which increases by~1
at each prime and is constant between consecutive primes.
The Mertens constant $\gamma_M$ governs the average drift
of this sawtooth via the asymptotic
\[
\sum_{p\le\kappa}\frac{\log p}{p}
= \log\kappa + O(1),
\]
connecting the renormalized scale to classical prime-sum estimates.

\begin{remark*}
The Mertens constant
\[
\gamma_M = \gamma + \sum_p\!\Bigl[\log\tfrac{p}{p-1}-\tfrac1p\Bigr]
\approx 0.2615,
\]
where $\gamma \approx 0.5772$ is the Euler--Mascheroni constant,
appears in prime-sum estimates that control the average behavior
of $H_{\mathrm{local}}^{\mathrm{ren}}$.
The same constant $\gamma$ appears in the first Li coefficient
$\lambda_1 = 1 + \gamma/2 - \log 2 - (\log\pi)/2 \approx 0.0231$,
reflecting the universal role of $\gamma$ in the transition
from discrete prime sums to continuous logarithmic integrals.
\end{remark*}

\begin{remark*}
The renormalized construction clarifies the scale of the prime-side
contribution, but does not by itself resolve the spectral inequality
\eqref{eq:openineq}.
\end{remark*}

\section{The remaining spectral inequality}\label{sec:open}

For the family $g^*_{1/2,\varepsilon}$, Weil positivity reduces to
\begin{equation}\label{eq:openineq}
Z(g^*_{1/2,\varepsilon})
\;\ge\;
H_{\mathrm{local}}(1/2,\kappa).
\end{equation}

\textbf{Open question.}
Is there an unconditional lower bound
\[
Z(g^*_{1/2,\kappa})
= \sum_{\rho}\bigl|\widehat g^*_{1/2,\kappa}(\rho)\bigr|^2
\;\ge\;
C\cdot(\log\kappa)^2
\]
for some constant $C>0$, where the sum runs over nontrivial zeros of~$\zeta$?

Formally, treating the zeros as independent via Selberg's
orthogonality principle~\cite{selberg} yields
$Z(g^*)\sim D\log\kappa\log\log\kappa$ with $D\approx 9.471$.
More precisely: the diagonal contribution
\[
\sum_p \bigl[c_p^{\mathrm{ren}}\bigr]^2 \cdot I(\varepsilon)
\;\sim\; D \cdot \log\kappa
\]
is unconditional (it follows from the prime number theorem alone);
the off-diagonal control of the cross-terms
$\sum_\gamma \cos(\gamma\log(p/q))$ for $p\neq q$
remains open. An unconditional proof of even
$Z\ge C(\log\kappa)^2$ is presently unavailable.

\begin{remark*}[The constant $D/2\pi$]
The asymptotic $Z(g^*)\sim D\log\kappa\log\log\kappa$ yields a
natural bridge constant
\[
\frac{D}{2\pi} \approx \frac{9.471}{6.283} \approx 1.507,
\]
which measures the ratio of the prime-side energy~$D$
to the archimedean normalization~$2\pi$ of the zero density
$N(T)\sim T\log T/(2\pi)$.
This constant is of a different nature from the pair-correlation
statistics studied in the GUE conjecture
(Montgomery~1973): it depends only on the total density of zeros
via the classical formula for~$N(T)$, which is unconditional
and requires neither GUE nor RH as input.
Whether $D/2\pi$ admits a representation in terms of
classical constants is an open question.
\end{remark*}

% FIX v5 — clarify status upfront
\begin{remark*}[Status of the spectral inequality]
The inequality~\eqref{eq:openineq} is open and of RH strength for this family.
By Lagarias~\cite{lagarias},
$\mathrm{RH} \iff W(f*\check f)\ge 0$ for all $f\in S_{ad}$.
A positive answer to the open question above would establish
$W(g^*_{1/2,\kappa}*\check g^*_{1/2,\kappa})\ge 0$ for the specific
family constructed here, which is a necessary but not sufficient
condition for RH in the absence of a density argument for this
family in~$S_{ad}$.
\end{remark*}

\section*{Computational note}

All numerical claims are reproducible. Code is available at~\cite{github}.
The Gaussian parameter $\varepsilon = 0.05$ is required for meaningful
numerical experiments with the zero sum (for $\varepsilon = 0.5$ one has
$e^{-\varepsilon^2\gamma_1^2} \approx 10^{-22}$, effectively
suppressing all zero contributions).

\medskip
\textbf{Numerical comparison $Z_{\mathrm{ren}}$ vs $H_{\mathrm{local}}^{\mathrm{ren}}$.}
With $c_p^{\mathrm{ren}} = f_p^{1/2}$, $\varepsilon=0.05$, and $N=100$ zeros:

\begin{center}
\begin{tabular}{ccccc}
\toprule
$\kappa$ & $Z_{\mathrm{ren}}$ & $H_{\mathrm{local}}^{\mathrm{ren}}$
         & $Z - H$ & $Z \ge H$ \\
\midrule
10   & 4.888 & 3.044 & $+1.844$ & $\checkmark$ \\
20   & 3.782 & 4.770 & $-0.988$ & $\times$ \\
29   & 5.520 & 3.588 & $+1.931$ & $\checkmark$ \\
50   & 5.693 & 2.881 & $+2.812$ & $\checkmark$ \\
100  & 5.311 & 1.562 & $+3.748$ & $\checkmark$ \\
200  & 5.102 & 2.355 & $+2.747$ & $\checkmark$ \\
500  & 5.156 & 1.101 & $+4.056$ & $\checkmark$ \\
1000 & 4.289 & 0.426 & $+3.863$ & $\checkmark$ \\
\bottomrule
\end{tabular}
\end{center}

\begin{remark*}
The table illustrates the sawtooth behavior of
$H_{\mathrm{local}}^{\mathrm{ren}}$: at each new prime $p$,
$H_{\mathrm{local}}^{\mathrm{ren}}$ jumps upward by $f_p > 0$,
then decreases continuously until the next prime.
This creates temporary inversions $Z_{\mathrm{ren}} < H_{\mathrm{local}}^{\mathrm{ren}}$
directly after a prime jump — exactly as $\pi(x) - \mathrm{Li}(x)$
fluctuates near small primes before stabilizing asymptotically.

In our experiments, violations occur for $\kappa \le 25$
(in the range between the primes 23 and 29).
For all $\kappa \ge 26$, $Z_{\mathrm{ren}} > H_{\mathrm{local}}^{\mathrm{ren}}$
holds consistently, with the ratio $Z/H$ growing from $1.5$ at $\kappa=29$
to over $10$ at $\kappa=1000$.
This asymptotic stability is consistent with the spectral
inequality~\eqref{eq:openineq}, though an analytic proof remains open.
\end{remark*}

\begin{thebibliography}{99}

\bibitem{paper1}
U.~Tehrani,
\emph{A Curvature Decomposition of the Explicit Formula
for the Riemann Zeta Function},
Research Note, March 2026, v6,
\href{https://doi.org/10.5281/zenodo.19025598}{%
doi:10.5281/zenodo.19025598}.

\bibitem{github}
U.~Tehrani,
\emph{AnalysisLab: Curvature and the Weil Functional},
GitHub repository, 2026,
\url{https://github.com/utehrani/analysislab-nt}.

\bibitem{lagarias}
J.~C.~Lagarias,
\emph{Li coefficients for automorphic L-functions},
Ann.\ Inst.\ Fourier \textbf{57}(5) (2007), 1689--1740.

\bibitem{li}
X.-J.~Li,
\emph{The positivity of a sequence of numbers
and the Riemann hypothesis},
J.\ Number Theory \textbf{65} (1997), 325--333.

\bibitem{weil}
A.~Weil,
\emph{Sur les ``formules explicites'' de la
th\'eorie des nombres premiers},
Comm.\ S\'em.\ Math.\ Univ.\ Lund (1952), 252--265.

\bibitem{mertens}
F.~Mertens,
\emph{Ein Beitrag zur analytischen Zahlentheorie},
J.\ Reine Angew.\ Math.\ \textbf{78} (1874), 46--62.

\bibitem{montgomery}
H.~L.~Montgomery,
\emph{The pair correlation of zeros of the zeta function},
in: Analytic Number Theory, Proc.\ Sympos.\ Pure Math.\ \textbf{24},
Amer.\ Math.\ Soc., Providence, 1973, pp.\ 181--193.

\bibitem{selberg}
A.~Selberg,
\emph{Contributions to the theory of the Riemann zeta-function},
Arch.\ Math.\ Naturvid.\ \textbf{48}(5) (1946), 89--155.

\bibitem{connes}
A.~Connes and C.~Consani,
\emph{Weil positivity and trace formula, the archimedean place},
Sel.\ Math.\ (N.S.) \textbf{27} (2021), Article~77.

\bibitem{edwards}
H.~M.~Edwards,
\emph{Riemann's Zeta Function},
Dover, 1974.

\end{thebibliography}

\section*{Acknowledgment}

The author thanks Jean-Fran\c{c}ois Burnol for confirming that the
curvature decomposition of~\cite{paper1} does not appear in the
literature known to him.

% FIX v5 — changelog
\section*{Change Log}
\textbf{v5 (April 2026):}
Framing of Weil positivity goal clarified;
spectral inequality explicitly marked as open and of RH strength.
All formal proofs, constructions and numerical results unchanged from v4.

\textbf{v4 (March 2026):}
Table values for $V_p(1/2)$ corrected; $c_p^{\mathrm{ren}}$ simplified
to $f_p^{1/2}$; $H_{\mathrm{ren}}$ asymptotic corrected.

\end{document}
